% NAF 9544 — lecture modernisée du volume entier.
%
% Pass: Opus 5.5 (claude-opus-5-5), 2026-09-22 — first pass, unchecked against
% the views by a human, and an interpretation rather than a record. Where this
% file and the transcriptions differ in substance, the transcriptions
% (batch-01.fr.tex, batch-02.fr.tex, batch-03.fr.tex) are the record.
%
% Sources read: the three batch transcriptions of this volume, in full, twice
% — once for the content, once for the apparatus — and, for the one
% concordance the description of view 54 invites, the project's own
% transcription of Latin 11195 (transcripts/latin-11195/batch-01.fr.tex and
% batch-02.fr.tex, header comments and views 1–3). No image was consulted. No
% printed edition was consulted to fill a gap or settle a word: not de Waard's
% Correspondance du P. Marin Mersenne, not the Œuvres philosophiques of Sophie
% Germain, not the Mémoires de l'Académie that Libri cites. Where the
% transcription has a gap, this reading says so. The three batches were
% transcribed on Opus 5.5 and cover views 1–57, the whole volume.
%
% What the volume turned out to be. An album of autographs bound in 1899
% (« Volume de 36 Feuillets — 27 Novembre 1899 »), UNSORTED: seven groups of
% pieces from 1620 to 1884, each announced by a dealer's cutting or a
% collector's wrapper, with nothing binding them but the sale room. In the
% order of the album:
%   v. 4–6     d'Alembert, minute to Malesherbes (annotated 1765).
%   v. 7–19    Bossut: letter to an unnamed count (12 Jan. 177[4]); letter to
%              Antide Janvier (2 Nov. an I); receipt to Agasse (3 fructidor
%              an X); one-leaf summary of his memoir on vaults (21 messidor
%              an IX), with two reading attestations.
%   v. 20–23   Sophie Germain, an untitled philosophical fragment on « l'idée
%              fondamentale d'une science ».
%   v. 24–28   Two telegrams of 28 and 30 June 1884 recording its purchase for
%              Baldassarre Boncompagni.
%   v. 30–36   Latin epigrams on Richelieu's death, attributed by the cutting
%              to Constantin Huygens; the leaves name no author.
%   v. 33–40   Mersenne: observations of 1620–1621 for the height of the pole
%              at Paris, addressed to Petau on the fold, and the letter to
%              Petau that sent them.
%   v. 41–55   G. Libri, « Description des manuscrits de Roberval », five
%              volumes, the last being the old Supplément latin 218, whose
%              tome I is today's Latin 11195.
% The résumé says the volume is unsorted; each group gets its own section.
%
% Mathematics. Thin, and not padded. Mersenne's spherical astronomy is redone
% in full: the conversion of ecliptic to equatorial coordinates, done on the
% leaf both by ruler and compass on a projection onto the solstitial colure
% and by sine tables, is shown to be exactly today's formula; every number was
% recomputed (declination right to one second; one subtraction off by 5'',
% which the transcription had noted; the polar-star mean and half-difference;
% the two solar declinations, right to 3'' with Tycho's obliquity). Bossut's
% vaults are a table of contents: modern names are attached with care and no
% calculation is supplied. Germain's fragment is epistemology of mathematics.
% Libri's statements that can be checked from the leaf alone were checked
% (the quadratures and the ratios of solids; the ellipse; the lemma of p. 720,
% which needs « arbitrarily little » to be true).
%
% Notation. The zodiac signs, set \add{Lion} etc. in the transcription, are
% given as ecliptic longitudes in degrees from the vernal point, with one
% footnote on the signs. The leaf's « 34.35' » style of angles becomes
% 34°35'. Sines with radius 100000 become sines with radius 1.
%
% Attributions stay as the leaves and guards give them. The Huygens
% attribution is the dealer's cutting; the Germain attribution rests on the
% heading and the wrapper; the Mersenne letter is signed « Mar. Mers. M. »,
% the observation sheet is not signed. The treatise on equations in Latin
% 11195 is not attributed here: Libri files it among Roberval's manuscripts,
% and this reading reports that as Libri's classing.
%
% Numbers on the leaves. The transcriptions write no \folio{}: a running ink
% series, one figure per leaf from 1 (v. 4) to 36 (v. 54), behaves as a
% foliation but was not taken as the library's. This reading cites views only.
\documentclass[11pt,a4paper]{article}
% The preamble shared by every transcription of the mathematicians' manuscripts in Gallica.
%
% It rests on one idea: **uncertainty compounds, it does not resolve.** A
% transcription that smooths over a doubtful word has destroyed the only thing
% it was made for — a reader must be able to tell, at every point, what was on
% the page and what was supplied. The apparatus macros below are therefore the
% critical apparatus, not decoration.
%
% The same commands are recognised by `scripts/render.mjs`, which produces the
% site's reading view, and by `scripts/tei.mjs`. A macro added here must be
% added there too, or rendering fails loudly — which is the wanted behaviour.
%
% Comments here are English, like the rest of the repository. The documents
% this preamble sets are French, and that is deliberate: see the skills.

\ProvidesPackage{germain}[2026/09/15 Transcriptions of mathematicians' manuscripts in Gallica]

\usepackage{amsmath,amssymb,amsthm}
\usepackage{mathrsfs}
% Kept from the parent preamble so the renderer's diagram support stays
% available; these manuscripts draw no commutative diagrams, and nothing here uses it.
\usepackage{tikz-cd}
\usepackage[english,french]{babel}
\usepackage{fontspec}

% --- Greek ------------------------------------------------------------------
% The text face stays fontspec's default, Latin Modern, which has no Greek:
% XeTeX drops every glyph it lacks with a line in the log and nothing on the
% page, and the Greek words the manuscripts quote vanished from the PDFs. So
% the Greek and Greek Extended blocks get a XeTeX character class of their own,
% and entering or leaving a run of it switches the family to CMU Serif — the
% same Computer Modern design, with polytonic Greek — and back. Only the
% family changes: a Greek word in \textit or \textbf keeps its shape and series.
% The transitions to and from every other class carry the tokens that class
% has towards an ordinary letter, so French spacing before « ; » or inside
% « » still holds after a Greek word. No grouping, so no brace or \endgroup
% can come out unbalanced; maths is left alone, since XeTeX inserts nothing
% there. Kept identical in `exercices.sty`.
\makeatletter
\newfontfamily\germainGreekfont{cmun}[
  Extension=.otf, NFSSFamily=germaingreek,
  UprightFont=*rm, ItalicFont=*ti, SlantedFont=*sl,
  BoldFont=*bx, BoldItalicFont=*bi, BoldSlantedFont=*bl]
\newXeTeXintercharclass\germain@greek
\def\germain@greekrange#1#2{%
  \@tempcnta=#1\relax
  \loop
    \XeTeXcharclass\@tempcnta=\germain@greek
    \ifnum\@tempcnta<#2\advance\@tempcnta\@ne\repeat}
\germain@greekrange{"0370}{"03FF}
\germain@greekrange{"1F00}{"1FFF}
\def\germain@greekprev{\rmdefault}
\def\germain@greekin{%
  \edef\germain@greekprev{\f@family}\fontfamily{germaingreek}\selectfont}
\def\germain@greekout{\fontfamily{\germain@greekprev}\selectfont}
\def\germain@greekjoin{%
  \XeTeXinterchartokenstate=\@ne
  \@tempcnta=\z@
  \loop
    \ifnum\@tempcnta=\germain@greek\else
      \XeTeXinterchartoks\germain@greek\@tempcnta=\expandafter{%
        \expandafter\germain@greekout\the\XeTeXinterchartoks\z@\@tempcnta}%
      \XeTeXinterchartoks\@tempcnta\germain@greek=\expandafter{%
        \the\XeTeXinterchartoks\@tempcnta\z@\germain@greekin}%
    \fi
    \ifnum\@tempcnta<4095\advance\@tempcnta\@ne\repeat}
% babel-french sets its own classes, and saves and restores the interchar
% state, each time a language is selected: join again after it does.
\addto\extrasfrench{\germain@greekjoin}
\addto\extrasenglish{\germain@greekjoin}
\AtBeginDocument{\germain@greekjoin}
\makeatother

\usepackage{geometry}
\geometry{a4paper,margin=25mm,marginparwidth=28mm}
\usepackage{marginnote}
\usepackage{xcolor}
\usepackage[normalem]{ulem}
\setlength{\emergencystretch}{3em}
\usepackage{hyperref}
\hypersetup{colorlinks=true,linkcolor=black,urlcolor=black,pdfborder={0 0 0}}

% --- Batch metadata ---------------------------------------------------------
% Taken as they stand by the rendering script, which shows them at the head of
% the reading view. They fix what the file claims to cover. The macro names are
% the parent project's, so that the scripts stay shared; \folder here means the
% volume's slug — `fr-9115` — and \pages counts Gallica views.

\newcommand{\folder}[1]{\gdef\grothFolder{#1}}
\newcommand{\batch}[1]{\gdef\grothBatch{#1}}
\newcommand{\pages}[2]{\gdef\grothFirst{#1}\gdef\grothLast{#2}}
\newcommand{\foldertitle}[1]{\gdef\grothFoldertitle{#1}}

% The holder's own shelfmark, printed as they write it: « Français 9115 ».
\newcommand{\shelfmark}[1]{\gdef\grothShelfmark{#1}}

% The dating, copied verbatim from the holder's catalogue — never inferred
% here. « XIXe siècle » is the BnF's; a pass that dates a leaf from its content
% says so in a \note{}, as its own inference.
\newcommand{\dating}[1]{\gdef\grothDating{#1}}

% The status notice, printed in the title block and repeated as a diagonal
% watermark on every page of the PDF. The manuscripts are in the public
% domain; what this line declares is not a right but a quality — an
% unverified first pass by a machine. The skills write the line; a file
% without it gets no watermark, so leaving it out is visible in the output.
\newcommand{\watermark}[1]{\gdef\grothWatermark{#1}}

\gdef\grothFolder{?}\gdef\grothBatch{}\gdef\grothFirst{?}\gdef\grothLast{?}%
\gdef\grothFoldertitle{}\gdef\grothShelfmark{}\gdef\grothDating{}\gdef\grothWatermark{}

% --- The résumé -------------------------------------------------------------
\newenvironment{resume}
  {\small\setlength{\parskip}{0.45em}}
  {\par\medskip\hrule\medskip}

% --- Keywords ---------------------------------------------------------------
% In English, closing the modernised reading's résumé: the modern vocabulary
% under which to search for what the leaves construct. The manifest extracts
% them to tag the volume — this line is the single source of the tags.
\newcommand{\keywords}[1]{\par\smallskip\noindent
  {\footnotesize\textsc{Keywords} --- \textit{#1}}\par}

% --- The critical apparatus -------------------------------------------------

% The Gallica view number, in the margin. This is the anchor that lets a
% disagreement over a formula be settled by looking at *one* image rather than
% a volume of 750. The site uses it to turn the facsimile as the transcription
% is scrolled. A view is one scanned image: usually one side of a leaf.
\newcommand{\page}[1]{%
  \marginnote{\footnotesize\color{gray!70}\textsf{v.\,#1}}%
  \label{page:#1}\ignorespaces}

% The BnF's pencilled foliation, where the leaf carries it and a pass can read
% it: « 198r ». This is what the literature cites — « ff. 198r–208v » — and
% the one line that turns a citation into a view. Recorded, never inferred:
% a folio a pass cannot read is not written.
\newcommand{\folio}[1]{%
  \marginnote{\footnotesize\color{gray!50}\textsf{f.\,#1}}\ignorespaces}

% The modernised reading's coarser counterpart to \page{}: which views a
% section is drawn from.
\newcommand{\pagerange}[2]{%
  \marginnote{\footnotesize\color{gray!70}\textsf{v.\,#1--#2}}%
  \ignorespaces}

% Illegible: never guessed.
\newcommand{\ill}{\textcolor{red!60!black}{[\,\ldots\,]}}

% A doubtful reading: the word is given, and flagged as uncertain.
\newcommand{\uncertain}[1]{\uline{#1}}

% An editorial addition — a missing letter, an implied word.
\newcommand{\add}[1]{\textcolor{blue!50!black}{[#1]}}

% Struck out by the writer. What was crossed out is part of the document.
\newcommand{\struck}[1]{\sout{#1}}

% The transcriber's note, on the state of the page or the mathematics.
\newcommand{\note}[1]{\par\smallskip{\small\color{gray!60!black}\textit{[#1]}}\par\smallskip}

% A marginal note of hers, distinct from the body of the text.
\newcommand{\marginal}[1]{\par\smallskip\noindent\hspace{1em}%
  {\small\color{gray!60!black}#1}\par\smallskip}

% --- Title ------------------------------------------------------------------
% The title block is French, like the documents themselves.

\AddToHook{shipout/background}{%
  \ifx\grothWatermark\empty\else
    \begin{tikzpicture}[remember picture, overlay]
      \node[rotate=52, scale=3.4, align=center, text=gray!18,
            font=\sffamily\bfseries]
        at (current page.center) {\grothWatermark};
    \end{tikzpicture}%
  \fi}

\AtBeginDocument{%
  \begin{center}
    {\large\bfseries Manuscrits de mathématiciens (Gallica) ---
      \ifx\grothShelfmark\empty\grothFolder\else\grothShelfmark\fi}\\[2pt]
    {\itshape\grothFoldertitle}\\[4pt]
    {\small vues \grothFirst--\grothLast
      \ifx\grothBatch\empty\else\ (lot \grothBatch)\fi}\\[2pt]
    \ifx\grothDating\empty\else
      {\small\color{gray!60!black}Datation du catalogue : \grothDating}\\[2pt]
    \fi
    \ifx\grothWatermark\empty\else
      {\small\bfseries\color{red!45!gray}\grothWatermark}\\[2pt]
    \fi
    {\footnotesize\color{gray!60!black}Transcription établie d'après les images IIIF de Gallica.
     Source gallica.bnf.fr / Bibliothèque nationale de France.\\
     Les lectures incertaines sont soulignées, les illisibles notées [\,\ldots\,],
     les ajouts entre crochets ; \textsf{v.} numérote les vues de Gallica,
     \textsf{f.} la foliotation au crayon de la BnF.}
  \end{center}
  \vspace{1em}}

\endinput


\folder{naf-9544}
\shelfmark{NAF 9544}
\pages{1}{57}
\dating{1601-1700}
\watermark{Lecture automatique\\interprétation non vérifiée}
\foldertitle{Recueil de lettres autographes de D'ALEMBERT, l'abbé BOSSUT, Sophie GERMAIN, C. HUYGHENS et le P. MERSENNE. — lecture modernisée du volume entier}

\begin{document}

\section*{Résumé}

\begin{resume}
Ce volume n'est pas un livre mais un album : trente-six feuillets reliés en
1899, sur lesquels on a monté, pièce à pièce, des autographes achetés en vente
publique, chacun précédé de la coupure du catalogue qui l'avait annoncé. Il
réunit, dans cet ordre, une minute de d'Alembert de 1765, quatre pièces de
l'abbé Bossut dont le sommaire d'un mémoire sur l'équilibre des voûtes, un
fragment philosophique de Sophie Germain suivi des deux télégrammes de 1884
qui en annoncent l'achat, des épigrammes latines sur la mort de Richelieu,
des observations astronomiques de 1620 et 1621 envoyées par Mersenne au
P. Petau avec la lettre qui les accompagnait, et enfin une description, par
Guglielmo Libri, de cinq manuscrits de Roberval. Il n'y a pas de fil entre
ces pièces, et cette lecture n'en invente pas : ce qui les rapproche est la
salle des ventes, non le sujet. La datation du catalogue, « 1601-1700 », ne
couvre que deux d'entre elles.

Les mathématiques y sont rares, et il faut les prendre où elles sont. Le
morceau le plus substantiel est la feuille de Mersenne. Le problème est celui
de la latitude de Paris, que l'on appelait la hauteur du pôle. On observe une
étoile voisine du pôle au moment où elle passe au méridien sous le pôle ; sa
hauteur au-dessus de l'horizon, augmentée de sa distance au pôle, donne la
hauteur du pôle. Encore faut-il connaître cette distance, c'est-à-dire la
déclinaison de l'étoile, alors que les catalogues donnent ses coordonnées par
rapport à l'écliptique. Mersenne fait la conversion de deux manières : à la
règle et au compas, sur une projection de la sphère sur le plan du méridien,
et par les tables de sinus. Sa construction est exactement la formule que
l'on écrit aujourd'hui ; son calcul est juste à la seconde près, sauf une
soustraction fausse de cinq secondes. Il y joint deux autres méthodes — la
moyenne des hauteurs de l'étoile polaire au-dessus et au-dessous du pôle, et
la hauteur du Soleil au solstice d'hiver — qui donnent trois valeurs entre
$48^\circ 41'$ et $48^\circ 48'$. Sa propre raison d'écrire tient en une
phrase de la lettre : il se souvient de ce que Petau lui avait demandé
« l'hiver passé sur la hauteur du pôle pour la latitude de Paris », il lui
envoie ses observations, et il le prie de les lui rendre, car il n'en a pas
gardé copie.

Les autres pièces ont chacune leur idée. Le sommaire de Bossut, lu en
séance le 10 juillet 1801, découpe la théorie des voûtes en berceau et en
dôme : épaisseur des pieds-droits qui résistent à la poussée, rupture en des
points donnés des reins, relation entre la forme de la voûte et les charges
qui pressent les voussoirs ; les calculs n'y sont pas. Le fragment de Sophie
Germain demande ce qu'est l'idée fondamentale d'une science : la définition
détermine son objet également pour tous, répond-elle, mais le fait connaître
inégalement ; elle refuse qu'un raisonnement ne soit qu'une opération
mécanique sur des signes et place la certitude mathématique dans l'identité
du sujet dont on parle. La description de Libri, enfin, est un inventaire de
Roberval : éléments de géométrie, quadrature des paraboles de tous les degrés
par les indivisibles, rapports des solides de révolution aux cylindres,
tangentes par la composition des mouvements, mécanique, hydraulique, lieux
géométriques, théorie des équations. Son dernier volume, le « supplément
latin 218 », se reconnaît aux marques qu'il cite : son tome I est
l'actuel Latin 11195.

Où cela mène. Du côté de Mersenne, à l'astronomie sphérique — coordonnées
écliptiques et équatoriales, déclinaison, détermination de la latitude par
les étoiles circumpolaires, réfraction atmosphérique, projection orthographique
de l'analemme. Du côté de Bossut, à la théorie des arcs et des voûtes en
maçonnerie, des mécanismes de rupture de Coulomb au calcul à la rupture. Du
côté de Germain, à la philosophie des mathématiques, à Condillac et à sa
« langue des calculs », et à ce qu'on appellera le formalisme. Du côté de
Libri, à Roberval : méthode des indivisibles, quadrature des paraboles
supérieures, méthode des tangentes par les mouvements composés, cycloïde et
sa compagne, centres de gravité, lieux géométriques, théorie des équations.

\keywords{spherical astronomy, ecliptic coordinates, declination, latitude determination, circumpolar stars, lower culmination, analemma, orthographic projection, astronomical refraction, theory of masonry arches, masonry domes, thrust, philosophy of mathematics, Condillac, formalism, method of indivisibles, method of exhaustion, quadrature of higher parabolas, solids of revolution, method of tangents by composition of motions, cycloid, centre of gravity, Archimedes' principle, geometric loci, theory of equations}
\end{resume}

\section*{Ce que le volume contient}

\pagerange{1}{57}
Cinquante-sept vues, dont les deux premières et les trois dernières sont la
reliure et ses gardes. La page de titre de l'album (vue 3) porte « Autographes
de Savants — Volume de 36 Feuillets — 27 Novembre 1899 ». Chaque pièce est
montée sur onglet ou sur un feuillet de montage et, le plus souvent, précédée
d'une coupure imprimée du catalogue de vente qui l'avait annoncée ; toutes
portent la marque « R. D. 9555 » et le cachet de la Bibliothèque
nationale.\footnote{Plusieurs séries de nombres courent sur les feuillets :
un chiffre à l'encre en haut à droite de chaque feuillet, de 1 (vue 4) à 36
(vue 54), sans lacune, qui s'accorde avec le compte de 36 feuillets du
certificat de 1899 et se comporte en tout comme une foliotation ; de grands
numéros de lot au crayon ou à l'encre (207, 96, 95, 97, 677, 335, 459) ; des
numéros de collections antérieures ; enfin la pagination propre de la
description de Libri et ses renvois marginaux. Les transcriptions n'ont
inscrit aucun folio, la première série n'ayant pu être tenue pour la
foliotation de la bibliothèque sur une numérisation en niveaux de gris, et
cette lecture renvoie donc aux vues. Les vues 13 et 50, secondes prises de vue d'une page déjà transcrite, et
les vues 56 et 57, la reliure, n'ont pas été transcrites, non plus que les
ouvertures blanches 29 et 34 ; les vues 26, 28 et 35 à 37 répètent des pages
déjà données.}

La datation du catalogue, « 1601-1700 », est reprise telle quelle en tête de
ce fichier comme des transcriptions. Elle ne vaut que pour les épigrammes et
pour les pièces de Mersenne ; le reste s'étend de 1765 à 1884. Le titre du
catalogue nomme cinq auteurs et en omet un, Libri, dont la description occupe
pourtant le tiers de l'album.

\section*{La minute de d'Alembert à Malesherbes (vues 4 à 6)}

\pagerange{4}{6}
Deux petits feuillets, de la main rapide de d'Alembert, écrits à la troisième
personne et sans signature : une minute, ou un mémoire destiné à être remis,
plutôt qu'une lettre.\footnote{La coupure de catalogue collée sur le feuillet
de montage (vue 4) décrit la pièce comme une « Let. aut. (minute) à
Malesherbes ; (1765), 2 p. 1/4 in-8 ». En tête, une main plus petite et plus
tardive a ajouté « a M. de Malesherbes » et « 1765 ». Le titre, souligné, est
lu « Lettre de M. d'alembert » avec doute : les lettres qui suivent « Lett »
se présentent autrement, et la lecture est celle du sens.} En voici le texte
en français d'aujourd'hui.

M. d'Alembert est pensionnaire surnuméraire dans la classe de géométrie depuis
neuf ans et demi. Il demande, et l'Académie demande pour lui, qu'il passe à la
place de pensionnaire ordinaire mécanicien. La seule objection qu'on puisse
lui faire, et qu'on lui fasse en effet, est la différence des classes :
objection sans fondement, puisqu'il y a mille exemples, à l'Académie, de
passages d'une classe à l'autre, et surtout de la géométrie à la mécanique et
réciproquement, ces deux classes ayant toujours roulé ensemble. M. de Montigny,
en 1758, est passé de la place d'associé géomètre à celle de pensionnaire
mécanicien ; M. d'Alembert demande beaucoup moins, puisqu'il est déjà
pensionnaire surnuméraire et qu'il ne s'agit que de le faire changer de
classe.

Cela ne porte préjudice à personne : d'abord parce qu'il est plus ancien que
tous les associés ; ensuite parce que M. de Courtivron, le plus ancien après
lui, demande la vétérance et a déclaré par lettre qu'il ne voulait point
concourir avec M. d'Alembert. Ainsi M. de Vaucanson, qui se trouve le second
associé mécanicien, deviendra le premier, et c'est tout ce qu'il pouvait
espérer, à supposer que les choses se passent à l'ordinaire et que M. de
Courtivron n'eût point demandé la vétérance.\footnote{« le second » est écrit
au-dessus de la ligne, au-dessus d'un mot raturé qui n'a pas été lu.}

M. de Saint-Florentin n'a donc autre chose à répondre à la demande de
l'Académie que oui ou non. M. d'Alembert espère que M. de Malesherbes voudra
bien représenter à M. de Saint-Florentin combien le non serait injuste,
malhonnête pour l'Académie, qui doit savoir ce qui lui convient le mieux, et,
on ose le dire, odieux après les travaux et les sacrifices de M.
d'Alembert.\footnote{« oui », « non » et le second « non » sont soulignés. Un
trait ferme le texte ; il n'y a ni date ni signature. Le comte de
Saint-Florentin était le ministre de la Maison du roi, dont relevait
l'Académie ; ce fait ne figure pas sur le feuillet. Le feuillet ne dit pas ce
qu'il advint de la demande.}

L'argument est de pure procédure : un changement de classe à grade constant
ne lèse aucun associé, puisque la place de mécanicien qui se libère par la
vétérance de Courtivron va de toute façon à Vaucanson. Rien ici n'est
mathématique, sinon la frontière — que d'Alembert tient pour poreuse — entre
la géométrie et la mécanique.

\section*{Les pièces de l'abbé Bossut (vues 7 à 19)}

\pagerange{7}{19}
Quatre pièces, chacune annoncée par un feuillet de montage, une fiche ou une
chemise : deux lettres, un reçu et le sommaire d'un mémoire.

\subsection*{La lettre à « Monsieur le Comte » (vues 7 à 10)}

Paris, 12 janvier 177[4].\footnote{Le dernier chiffre du millésime est ouvert
et se lit 4 plutôt que 6, sans certitude. Sous la date, d'une autre main et
d'une autre plume, sans doute celle du destinataire : « Réponse donnée ».}
Bossut s'excuse d'avoir tardé à répondre à la lettre du 29 novembre
précédent. L'honneur que la « célèbre académie » du comte lui a fait de
l'adopter le pénètre de reconnaissance ; il est glorieux de voir son nom
associé à ceux de tant de grands hommes ; il prie le comte de le mettre aux
pieds de l'académie et de lui témoigner sa vénération. Il s'est acquitté de la
commission du comte auprès de d'Alembert, qui le charge d'un million de
remerciements. Le comte a emporté l'estime et les regrets de tous ceux qui ont
eu le bonheur de vivre « ici » avec lui ; Bossut regrette de n'avoir fait que
l'entrevoir.\footnote{Le mot rendu par « tenir » (« je veux tenir mon pardon
uniquement de votre bonté ») est donné avec doute ; « m'adopter » est écrit
vite et lu par le sens ; « chagin » est écrit pour « chagrin ». La main de
Bossut forme l'r final en petite boucle et écrit « Conte » pour « Comte ».}
Ni le comte ni son académie ne sont nommés sur les feuillets, et cette lecture
ne les identifie pas.

\subsection*{La lettre à Antide Janvier (vues 11 à 16)}

Paris, 2 novembre de l'an I de la République, soit le 2 novembre
1792.\footnote{Le feuillet écrit « 2 9\textsuperscript{bre} l'an 1. de la
Rep. » : le mois est celui du calendrier grégorien, l'an celui de la
République, commencée le 22 septembre 1792. La coupure de catalogue de la
vue 11 donne la même date. Au-dessus de la coupure, au crayon, « cat. Dufour
30 janv. 1889 », le nom lu avec doute ; sur la fiche de la vue 12, au crayon,
« Arrigoni, XIII\textsuperscript{e} catalogue 1878 », l'année lue avec doute.}
Bossut remercie Janvier du portrait gravé de Lalande. Il refuse de le
conseiller sur la pétition que Janvier veut présenter à l'Assemblée nationale
— il l'avait même prié de ne pas lui en envoyer la minute. Mais il estime ses
talents : Janvier mérite un logement dans les maisons nationales, et Bossut
l'engage à demander au ministre celui que M. Paris occupait aux Menus, en le
priant de se faire rendre compte de ses talents par la commission du Muséum ;
Bossut, alors légalement en mesure de le servir, y mettra tout son
zèle.\footnote{Le nom « Paris » se lit aussi « Pavis » et reste douteux. La
phrase laisse entendre que Bossut siégeait à la commission du Muséum, sans le
dire.} Si l'affaire prend un autre chemin, il aura peine à réussir pour un
autre objet, les logements des Menus ne regardant pas en général la
commission ; mais il pourrait se faire que M. Roland désire connaître l'avis
de la commission sur Janvier.\footnote{Une tache couvre, dans la marge, la
fin de « il » et le mot avant « façon », lu « la » par le sens. Roland est
vraisemblablement Jean-Marie Roland, alors ministre de l'Intérieur ; le
feuillet ne donne que le nom.} L'adresse (vue 16) est : « À Monsieur Janvier,
horloger mécanicien, aux Menus, rue Poissonnière et Bergère, à Paris ».

\subsection*{Le reçu à Agasse (vue 17)}

« J'ai reçu du citoyen Agasse les 65\textsuperscript{e}, 66\textsuperscript{e}
et 67\textsuperscript{e} livraisons de l'Encyclopédie méthodique. À Paris,
ce 3 fructidor an X » — le 21 août 1802 — signé Bossut.\footnote{La
conversion est de cette lecture. Le chiffre du jour, un 3 bouclé, pourrait
être un 5, ce qui donnerait le 23 août 1802. Sous le billet est montée une
chemise d'une main italienne, qui porte trois notices : Bossut, Laplace et
Delambre, avec leurs dates de naissance et de mort ; les deux dernières
annoncent les attestations de la pièce suivante.}

\subsection*{Le sommaire du mémoire sur l'équilibre des voûtes (vues 18 et 19)}

Un feuillet daté du 21 messidor an IX, soit le 10 juillet 1801, de la main
de Bossut selon une note de collectionneur, avec dans la marge deux
attestations de lecture : « Lu le 21 messidor an 9 à la réserve des
calculs », d'une autre main, et « Le 21 messidor an 9 le cit. Bossut a lu le
mémoire dont cette feuille est le sommaire », signé Delambre.\footnote{Une
note moderne au crayon attribue la première attestation à Laplace
(« aut. de Laplace ») ; rien d'autre sur le feuillet ne l'établit. La
conversion de la date est de cette lecture.} Le texte :

Le mémoire est divisé en trois sections : la première traite de l'équilibre
des voûtes en berceau, la seconde de l'équilibre des voûtes en dôme, la
troisième contient des expériences et des observations sur la résistance des
pierres et sur la manière dont les voûtes trop chargées se
rompent.\footnote{Le feuillet a d'abord écrit, puis raturé, un ou deux mots
avant « rompent », dont l'un paraît être « voutes », avant de récrire « voutes »
au-dessus de la ligne entre « les » et « trop ». La phrase de la dernière
rédaction dit « la manière dont les voûtes rompent les voûtes trop
chargées » ; la lecture ci-dessus garde le sens et non la lettre.}

\begin{itemize}
  \item[I.] \emph{Voûtes en berceau.} Chapitre 1 : épaisseur des pieds-droits
  lorsque la voûte se rompt en des points donnés des reins. Chapitre 2 :
  relations générales entre la forme de la voûte et les forces qui pressent
  les voussoirs. Chapitre 3 : forme à donner aux faces latérales et
  extérieures des pieds-droits lorsqu'ils peuvent se rompre par tranches
  horizontales.
  \item[II.] \emph{Voûtes en dôme.} Chapitre 1 : épaisseur des pieds-droits
  pour résister à la poussée de la partie supérieure, formant une calotte que
  l'on considère comme un corps continu. Chapitre 2 : conditions de
  l'équilibre entre toutes les parties d'une voûte en dôme. Chapitre 3 : forme
  à donner aux faces latérales et extérieures des pieds-droits lorsqu'ils
  peuvent se diviser par tranches horizontales.
  \item[III.] Deux chapitres, l'un sur les expériences, l'autre sur les
  observations.
\end{itemize}

Le sommaire n'a pas de calcul, et l'attestation dit que les calculs n'ont
pas été lus en séance ; cette lecture n'en supplée aucun.\footnote{« se rompt
en » (section I, chapitre 1) est une lecture douteuse.} Ce qu'on peut dire,
c'est sous quels noms ces questions sont connues aujourd'hui. La voûte en
berceau est un arc prolongé ; les pieds-droits sont ses appuis ; les reins,
la partie de l'arc entre la clé et les naissances. Le premier chapitre suppose
la rupture en des points donnés, c'est-à-dire un mécanisme de ruine où la
voûte se divise en blocs qui tournent autour de quelques joints : c'est
l'approche de Coulomb (1773), qui cherche parmi ces mécanismes celui qui
exige le plus grand appui, et que le calcul à la rupture du
XX\textsuperscript{e} siècle a mise en forme de
théorème.\footnote{Le feuillet ne nomme personne ; le rapprochement est de
cette lecture et ne dit rien de ce que le mémoire contient. Coulomb est
antérieur au mémoire et Bossut pouvait le connaître ; la formulation moderne
par la ligne des pressions et par l'analyse limite est postérieure.} Le
deuxième chapitre est le problème inverse : la forme d'un arc en équilibre
sous des charges données, dont la chaînette renversée est le cas classique
d'un arc sous son seul poids. Le troisième chapitre de chaque section cherche
un profil des pieds-droits tel que chaque tranche horizontale soit en
équilibre, un problème de solide d'égale stabilité. La seconde section passe
à la symétrie de révolution, où la calotte traitée comme un corps continu
annonce une coque plutôt qu'un assemblage de voussoirs.

\section*{Sophie Germain : l'idée fondamentale d'une science (vues 20 à 23)}

\pagerange{20}{23}
Deux petits feuillets, paginés 16 et 17, sans titre, sans date ni signature.
Une ligne d'une autre plume les ouvre : « Ce qui suit s'est trouvé dans les
papiers laissés par Mlle Germain ». L'attribution repose sur cette ligne et
sur la chemise de la vue 20, « Sofia Germain (autogr.) », non sur un examen
de la main par cette lecture.\footnote{La transcription note que la main a
été comparée à un autographe de Sophie Germain (Français 9115, vue 222) et
qu'elle est compatible, sans aller plus loin. Plusieurs mots sont d'une encre
plus pâle que le reste de la ligne, comme s'ils avaient été écrits dans des
blancs laissés d'abord.} En voici le texte en français d'aujourd'hui.

\subsection*{Le texte}

Pour celui qui commence l'étude d'une science, l'idée fondamentale de cette
science n'est que ce qui en détermine le sujet. À mesure qu'il avance, il
apprend à connaître les divers points de vue sous lesquels l'objet de la
science peut être considéré ; et lorsqu'il est parvenu à la connaissance
parfaite du sujet, l'idée fondamentale est pour lui la science tout entière,
parce que cette idée s'est alors agrandie de toutes celles qu'il a acquises
depuis, et de leurs rapports.

Mais, a-t-on dit, si la science tout entière n'était pas contenue dans l'idée
fondamentale, comment pourrait-on l'y voir ?\footnote{Le mot « on » est
chargé d'encre ; il n'est pas sûr qu'il soit biffé plutôt que repassé.}

Je réponds : la proposition fondamentale d'une science, bien qu'exprimée en
termes exacts et constants, ne contient pas, pour tous ceux qui l'entendent,
une idée identiquement la même, mais seulement des idées qui s'accordent dans
leurs traits caractéristiques ; de sorte que, pour tous les esprits, le sujet
est également déterminé mais inégalement connu, et que cette proposition, qui
n'a d'abord été, à l'entrée de la carrière, qu'une idée caractéristique,
devient, lorsqu'on l'a parcourue, une véritable description du sujet, et
remplit au plus haut degré la fonction de définition par rapport à la science
dont elle détermine l'objet.\footnote{« à l'entrée de la carrière » est
ajouté au-dessus de la ligne, d'une plume plus fine, et se lit mal ; la
lecture en est douteuse.} En effet, le souvenir des propositions qui ont été
déduites de la proposition fondamentale, et dont la réunion forme l'ensemble
de la science, se lie, pour celui qui la possède, à l'énoncé de l'idée
fondamentale, de sorte qu'il voit d'un seul coup d'œil le sujet tout entier.

On a confondu l'idée fondamentale avec le sujet de la science ; c'est ainsi
que l'on a pu dire : \emph{l'inconnu est dans le connu ; lorsqu'on apprend,
on va du même au même ; une science ne renferme qu'une seule idée}. Rien
n'est plus absurde que de dire que l'on puisse avoir une idée sans la
connaître, et c'est pourtant ce que cette doctrine suppose.\footnote{Les
trois propositions citées sont soulignées sur le feuillet. « dire » remplace
« supposer », biffé.}

On est allé jusqu'à prétendre que, dans une suite de déductions, vers la
deuxième ou la troisième il ne restait plus d'idées mais seulement des signes
d'idées, puis des signes de signes, et qu'un raisonnement bien fait n'était
qu'une opération mécanique exécutée sur des signes, dont il fallait, pour
revenir aux idées, retourner jusqu'à l'idée fondamentale. Et, a-t-on dit, si
vous croyez avoir des idées dans une longue suite de déductions, comment ne
vous satisfaites-vous de leur démonstration qu'en revenant à l'idée
première ?\footnote{« prétendre » remplace « dire », biffé ; « puis » est
biffé avant « et qu'un raisonnement » ; au bas de la page de gauche,
« arriver » est barré avec deux ou trois lettres que le même trait paraît
traverser, et la phrase reprend en haut de la page de droite par « retourner
jusqu'à ». La lecture « comment ne vous satisfaites-vous » modernise « comment
ne vous contentez vous sur leur demonstration ».}

Je crois que l'on peut répondre : les propositions sont essentiellement des
idées toutes les fois qu'elles ont un sens intelligible. Si, dans une suite de
déductions, on sent le besoin de retourner en arrière, c'est uniquement pour
s'assurer qu'elles conviennent au sujet auquel on les attribue ; c'est enfin
pour s'assurer de l'identité du sujet énoncé avec le sujet qualifié. La
certitude mathématique tient à l'identité du sujet, et la langue des calculs
est bien faite parce qu'en raisonnant suivant les règles qu'elle prescrit, il
est impossible d'attribuer au sujet de la science des propriétés qui ne
soient pas contenues dans l'idée fondamentale ; de sorte que les questions
que l'on examine ne se trouvent jamais compliquées de considérations qui lui
soient étrangères.\footnote{Le fragment s'arrête ici, deux lignes en haut du
verso du second feuillet (vue 23), sans signature ni date. Rien ne le suit.}

\subsection*{Ce que le fragment soutient}

La doctrine visée n'est pas nommée. Ses trois formules — l'inconnu dans le
connu, le passage du même au même, la science réduite à une seule idée — et
l'expression « la langue des calculs » sont celles de Condillac et de l'école
des idéologues, pour qui toute évidence de raison est une identité et le
calcul une langue bien faite.\footnote{L'identification est de cette lecture ;
le feuillet ne nomme ni Condillac ni personne. \emph{La Langue des calculs}
est le titre de l'ouvrage posthume de Condillac (1798).} Sophie Germain
retourne la formule contre la doctrine. Elle accepte que la langue des calculs
soit bien faite, mais pour une autre raison : non parce qu'un raisonnement ne
serait qu'une suite d'identités vides, mais parce que ses règles empêchent
d'attribuer au sujet ce qui ne lui appartient pas. Elle distingue deux choses
que la doctrine confond : ce qui \emph{détermine} un objet, et ce qu'on en
\emph{connaît}. Une définition détermine l'objet pour tous ; la connaissance
des propriétés qui en découlent est l'œuvre du temps, et elle est
inégale.

En termes d'aujourd'hui, la distinction est celle d'un objet fixé par une
définition et de la théorie que l'on déduit de cette définition — théorie qui
n'est pas « connue » du seul fait qu'elle est déterminée. Et la position
qu'elle combat, celle d'un raisonnement réduit à une manipulation mécanique de
signes, est ce qu'on appellera plus tard le formalisme.\footnote{Le nom est
postérieur au feuillet et n'y figure pas ; il est donné ici pour dire sous
quel mot la question a été reprise, non pour prêter à Sophie Germain une
position dans un débat qui n'existait pas encore.} Sa réponse à l'objection —
on revient au principe pour vérifier que l'on parle toujours du même sujet —
fait de l'identité de référence, et non de la forme des signes, la garantie de
la déduction. Le fragment se lit comme une note préparatoire ; il n'est pas
daté, et le seul terme qu'il porte est celui que lui donne sa première ligne :
il était dans ses papiers à sa mort.

\section*{L'achat de 1884 (vues 20 et 24 à 28)}

\pagerange{20}{28}
Le fragment de Germain est précédé d'une chemise, « Sofia Germain (autogr.) »,
d'un grand numéro « 677 », et d'un billet : « Sous ce pli l'autographe de
Sophie Germain, acheté le lundi 30 juin (cabinet Dubrunfaut). La couverture
imprimée Charavay sera envoyée avec une brochure quelconque. A. M. »\footnote{Les
initiales « A. M. » ne sont pas autrement expliquées, et cette lecture ne les
résout pas. L'année n'est pas écrite. Sur la chemise, « CLXXI » et « 2810 » ;
sur la chemise et le billet, 65 sur 1 et 65 sur 2, et la même fraction, 65
sur 6 ou sur 7, reparaît au crayon sur les feuillets de Germain et sur les
télégrammes.} Il est suivi de deux télégrammes reçus à Rome, sur les
formules imprimées de l'administration italienne, adressés à « Balthasar
Boncompagni, Palais Piombino, Rome ».\footnote{L'une des bandes porte
« BONCOMPAGNIE » ; l'autre « BONCOMPAGNI ». Les vues 26 et 28 sont de
secondes prises de vue des vues 24 et 27.}

Le premier, dans l'ordre des dates (vues 27 et 28), est du 28 juin 1884 :
« Reçu télégramme achat autographe Germain ». Le second (vues 24 à 26),
déposé à Paris le 30 et reçu le 30 juin 1884 à 19 heures : « Germain acheté
29 part demain ».\footnote{Le dernier chiffre de l'année, ajouté à la main
sur la formule du premier télégramme, est mal formé ; la bande imprimée porte
« 28/6/84 ». Selon l'avis imprimé de la formule, les nombres de la ligne
d'en-tête donnent le numéro de la dépêche, le nombre de mots, la date et
l'heure du dépôt. Le « 29 » du second télégramme n'a pas d'unité : ce peut
être un prix, un numéro de lot ou autre chose, et cette lecture ne le décide
pas.} L'ordre des opérations se lit ainsi : Boncompagni, à Rome, télégraphie
un ordre d'achat ; son correspondant de Paris en accuse réception le 28 juin ;
l'achat est fait, et le 30 le même correspondant annonce que la pièce part
le lendemain. Le 30 juin 1884 était un lundi, ce qui accorde le billet de la
vue 20, qui ne porte pas d'année, avec les télégrammes.\footnote{Le calcul du
jour de la semaine est de cette lecture. Dubrunfaut est le nom du cabinet
d'autographes dispersé ; le billet ne dit rien de plus. La description de
Libri, plus loin dans l'album, porte au pied de sa première page « Vente
B. Boncompagni, VI. Autogr. (1898), n\textsuperscript{o} 673 » : les deux
pièces sont passées par la même collection, et le 677 de la chemise de
Germain pourrait être son numéro dans la même vente, ce que rien sur le
feuillet ne dit.}

\section*{Épigrammes latines sur la mort de Richelieu (vues 30 à 36)}

\pagerange{30}{36}
Deux feuillets d'une belle main calligraphique du XVII\textsuperscript{e}
siècle, paginés 21 et 22, portant huit pièces de vers latins. Aucun nom
d'auteur n'y figure. La coupure de catalogue de la garde (vue 30) les
attribue à Constantin Huygens, le poète et homme d'État hollandais, et les
résume : épigrammes sur le retour de Richelieu, sur sa mort, sur sa
bibliothèque, d'autres contre Louis XIII, la dernière adressée à Marie de
Médicis.\footnote{Sous la coupure : « Vendita Grimm. Romanzoff. Paris, mai
1886 », le second nom lu avec doute. Les vues 35 et 36 sont de secondes prises
de vue des vues 31 et 32. La main écrit l'e final et le second élément de æ
comme une petite boucle surélevée, l'o final comme un s arrondi, le t final
comme un s à longue queue, et une même capitale bouclée pour B, D, V et H ;
beaucoup de lectures se règlent sur le sens.} Ce ne sont pas des
mathématiques ; on en donne ici une traduction française en prose, les
lectures douteuses de la transcription étant signalées.

\emph{Sur le retour du cardinal Richelieu.} L'Espagnol pâlit, déchu d'un noir
espoir : le Gond sur lequel tourne la vaste machine du monde agit par un
prodige. L'ancienne vigueur est revenue aux nerfs de Richelieu ; le fléau
brûlant a quitté ses épaules d'Atlas. Acquitte ton vœu envers le chef qui
revient, ô France, afin qu'il continue de vaincre pour toi et de
t'embrasser.\footnote{\emph{Cardo}, le gond, joue sur \emph{Cardinalis}.
« prodigio » et le dernier mot, « habet », sont des lectures douteuses ; la
fin du distique est traduite par le sens général.}

\emph{Sur les détracteurs de Richelieu mort.} La foule qui, suppliante,
soumettait au cardinal vivant son esprit et sa main, et ne reconnaissait sur
terre aucune divinité plus présente, se rue sur l'homme enseveli avec des
sarcasmes, des plaisanteries, des injures, et vomit, impuissante, tout ce
qu'une bouche mordante peut cracher sur les mânes de celui qui ne le méritait
pas. Peuple léger, plèbe ingrate, inhumaine comme un serpent : l'Atlas,
l'Atride, le Gond du roi et du royaume, une fois né pour la patrie, n'aurait
pas dû cesser de naître ; mais pour toi, il n'aurait dû être ni mort ni
vivant.\footnote{Le titre porte « detractorem », au singulier, lu avec
doute ; le poème parle d'une foule. « Ruit », au quatrième vers, est offert
par le sens. « simel » est écrit ainsi pour \emph{semel}.}

\emph{Sur ce vers de Bourbon}, dont le titre cite le texte auquel il répond
(« … pâture des vers qui naissent »). On arrache indignement la barbe du lion
abattu ; la plèbe méchante frappe Richelieu de la dent ; elle l'emporte, mais
sur un mort : la descendance qui naguère adorait le vivant, suppliante, mord
à belles dents son père. Voilà — c'est la seule chose où tu ne te trompes
pas, impie — il gît, pâture des vers ingrats nés de lui-même.\footnote{La
seconde ligne du titre, abrégée (« Ext\textsuperscript{o} »), n'est pas
résolue, et son dernier mot est douteux ; « Vincit » et « patrem » sont des
lectures douteuses. Le feuillet ne dit pas quel Bourbon est cité.}

\emph{Sur la bibliothèque de Richelieu.} Mille livres, et mille livres, et
des milliers de livres, rassemblés par le soin public du magnifique Richelieu,
pour le bien de la patrie, dans de superbes armoires, dressent les marbres
vivants d'un nom éternel. Et pour que cette masse auguste de volumes achetés
grandisse toujours vers ses titres, il ordonna que des milliers de livres
fussent remplis de mille livres. Cesse, postérité, […], par où la masse, en
roulant, se double comme une boule de neige. Il a laissé à vous qui écrivez
une matière à la mesure de vos forces, où toute la France peut suer. Écrivez
Richelieu, les triomphes remportés pour le roi, écrivez des actions à lire en
caractères éternels : la bibliothèque qui recevra la gloire posthume d'un seul
homme remplira mille livres.\footnote{« coemptis » (achetés) est lu pour
« coimptis » ; après « posteritas » un mot n'a pas été lu et le suivant est
biffé ; « habet » est de nouveau douteux. La phrase « il ordonna que des
milliers de livres fussent remplis de mille livres » traduit la lettre du
vers, dont le sens n'est pas clair.}

\emph{Sur sa mort.}\footnote{En tête de la page, en petit : « Audi alteram
partem », écoute l'autre partie.} Il convenait ainsi : Armand est parti en
compagnon de la dame étrusque. Français, pourquoi gémis-tu, dégénéré, d'une
bouche insensée ? Veuille être ce que tu es : te voilà enfin monarchie sous
un roi. Le Fils a succédé à la mère et au serviteur.\footnote{« Hetruscæ » est
une lecture douteuse. La dame étrusque, c'est-à-dire toscane, est Marie de
Médicis, morte quelques mois avant Richelieu ; le Fils est Louis XIII, qui
succède à sa mère et à son ministre, le « serviteur ». Ces identifications
sont de cette lecture.}

\emph{Sur le roi.} Que celui qui a perdu son bâton prenne garde que le Juste
ne tombe désormais plus de sept fois par jour.\footnote{\emph{Scipio}, le
bâton, joue sur le nom du général romain ; le Juste est le surnom de
Louis XIII, et le vers parodie le proverbe « le juste tombe sept fois ».}
Et : jusqu'ici l'Espagnol respectait non sans raison un roi bien armé — un
Armand ; il espère maintenant, non sans raison, qu'il sera désarmé.\footnote{Le
jeu porte sur \emph{Armandum} et \emph{exarmandum}. Une esperluette isolée
sépare ce distique de la pièce précédente.}

\emph{[Sans titre.]} Marie, mère d'une grâce impuissante, où que tu sièges
parmi les mânes, ayant obtenu du même coup le repos refusé à la vivante,
garde-toi du Cardinal, si tu es sage : il est suspect, celui qui a voulu
suivre si vite la morte qu'il avait persécutée vivante.\footnote{« ut »
(ou « et »), « Quiete » et le début de « consides », sous une tache, sont
douteux.}

Les pièces se situent d'elles-mêmes : après la mort de Richelieu, qui suit
celle de Marie de Médicis, et du vivant de Louis XIII, à qui l'une d'elles
s'adresse. C'est une datation par le contenu, de cette lecture, entre la fin
de 1642 et le printemps de 1643 ; aucune date n'est écrite.

\section*{Mersenne : les observations de 1620 et 1621 pour la hauteur du pôle à Paris (vues 33 à 39)}

\pagerange{33}{39}
La garde de la vue 33 porte la coupure d'une vente Charavay du 6 février
1889 : « 1° L. a. s. au Père Petau, 1 p. in-4 oblong. Belle lettre où il lui
annonce la résolution d'un problème relatif à la hauteur du pôle au-dessus de
l'horizon de Paris. 2° Pièce autographe, 1 p. 1/2 in-fol. C'est la résolution
du problème annoncée dans la lettre précédente. » La pièce est un feuillet
in-folio en latin (vue 38, et son verso à gauche de la vue 39), sans
signature ; plié en lettre, il porte au dos l'adresse « Au Révérend Père Denis
Petau, de la Société de Jésus, au Collège de Clermont ».\footnote{La vue 37
répète la vue 33. L'attribution à Mersenne est celle de la coupure et de la
lettre de la vue 39, qui annonce ces observations ; le feuillet lui-même ne
porte pas de nom.}

\subsection*{Le problème}

« L'an 1621, le 3 janvier, la brillante du quadrilatère de la Petite Ourse a
été vue au-dessus de l'horizon, dans le méridien, sous le pôle, à
$34^\circ 35'$ de hauteur. Au même moment, la longitude de l'étoile était de
$127^\circ 34' 20''$ comptés depuis l'équinoxe de printemps, sa latitude de
$72^\circ 51' 30''$. On cherche la déclinaison de l'étoile. »\footnote{Le
feuillet écrit la longitude « 7.34$'$.20$''$ » suivie du signe du Lion,
c'est-à-dire $7^\circ 34' 20''$ dans le signe du Lion, et donne aussitôt
l'équivalent « ab æquinoctio verno 127.34$'$.20$''$ ». L'usage du temps
comptait les longitudes écliptiques par signes de $30^\circ$ à partir du point
vernal — Bélier de $0^\circ$ à $30^\circ$, Taureau, Gémeaux, Cancer de
$90^\circ$ à $120^\circ$, Lion de $120^\circ$ à $150^\circ$, et ainsi de suite
jusqu'aux Poissons ; le point vernal est le « début du Bélier », le solstice
d'été le « début du Cancer ». Le feuillet dessine les signes ; la
transcription les rend par leur nom, et cette lecture par la longitude en
degrés de $0^\circ$ à $360^\circ$. Les angles sont écrits sur le feuillet avec
des points, « partium 34.35$'$ » ; ils sont notés ici à la manière d'aujourd'hui.
« stilo » est suivi d'un point, sans dire si le style est l'ancien ou le
nouveau.} L'étoile est aujourd'hui $\beta$ de la Petite Ourse, Kochab, la plus
brillante des quatre étoiles du quadrilatère.\footnote{L'identification est de
cette lecture, d'après la description du feuillet et d'après les coordonnées
qu'il donne, qui s'accordent avec la position de cette étoile au début du
XVII\textsuperscript{e} siècle.} La déclinaison obtenue, on en déduit
la hauteur du pôle.

Notons $\lambda$ la longitude écliptique, $\beta$ la latitude écliptique,
$\delta$ la déclinaison, $\varepsilon$ l'obliquité de l'écliptique et $\varphi$
la hauteur du pôle, c'est-à-dire la latitude du lieu. Le feuillet prend
$\varepsilon = 23^\circ 31' 30''$ « selon Tycho ».\footnote{Le nom se lit
« Tichonem » ou « Tirhonem » ; « Tychonem » est offert avec doute. Un trait en
boucle, peut-être un mot biffé, précède « maxima ».} Les données sont donc
\[
\lambda = 127^\circ 34' 20'', \qquad \beta = 72^\circ 51' 30'', \qquad
\varepsilon = 23^\circ 31' 30'', \qquad h = 34^\circ 35',
\]
où $h$ est la hauteur observée au passage inférieur.

\subsection*{La figure}

Le grand cercle de la figure, que le feuillet appelle le méridien, est le plan
qui contient à la fois l'axe du monde et l'axe de l'écliptique : le colure des
solstices. Son diamètre $\alpha\beta$ est la trace de l'équateur, $\gamma\delta$
l'axe du monde ($\gamma$ pôle nord, $\delta$ pôle sud), $\varepsilon\zeta$ la
trace de l'écliptique, $\eta$ le pôle nord de l'écliptique ; les arcs
$\alpha\varepsilon$, $\beta\zeta$, $\gamma\eta$ et $\delta\theta$ valent
l'obliquité. Le parallèle de latitude de l'étoile — le petit cercle des points
de latitude écliptique $\beta$ — se projette sur ce plan suivant la corde
$A\kappa$, avec $\varepsilon A = \zeta\kappa = 72^\circ 51' 30''$. Le point $A$
est le point de ce parallèle qui a la longitude $90^\circ$ (au-dessus du
solstice d'été), $\kappa$ celui qui a la longitude $270^\circ$. Rabattu dans le
plan de la figure, le petit cercle porte l'étoile en $\mu$, à l'arc
$A\mu = \lambda - 90^\circ = 37^\circ 34' 20''$ de $A$.\footnote{Le feuillet
écrit que $A\mu$ est la longitude « comptée depuis le point équinoxial ou
solsticial $A$ », 127°34'20'' ; puis « depuis ledit point solsticial,
$7^\circ 34' 20''$ du Lion, ou $37^\circ 34' 20''$ depuis le Cancer ». Les deux
nombres diffèrent de $90^\circ$ : $127^\circ 34' 20''$ est compté depuis le
point vernal, $37^\circ 34' 20''$ depuis le solstice, et c'est ce second arc
que le calcul emploie. La transcription donne la phrase telle qu'elle se lit.
La main écrit la lettre $A$ comme un $\lambda$ capital, distinct du $\lambda$
minuscule de la figure.}

\emph{La solution à la règle et au compas.} Tout étant tracé à l'échelle, on
abaisse de $\mu$ la perpendiculaire $\mu F$ sur la corde $A\kappa$ : $F$ est la
projection de l'étoile sur le plan du colure. Par $F$ on mène la parallèle à
l'équateur, qui coupe le cercle en $o$ et $\xi$ ; l'arc $\alpha o$ (ou
$\beta\xi$) est la déclinaison. C'est une projection orthographique de la
sphère sur un plan méridien, le procédé que la tradition appelle l'analemme :
toute la trigonométrie sphérique du problème est remplacée par une
construction plane.\footnote{Le nom d'analemme n'est pas sur le feuillet ; il
est donné ici parce que la construction coïncide avec celle que la tradition
grecque désigne ainsi. Le titre de la section du feuillet est « Mechanicè per
regulam et circinum », le « per » étant une lecture douteuse.}

\subsection*{Le calcul par les sinus}

Le feuillet refait la construction en nombres, avec des sinus de rayon
$100\,000$ ; on les écrit ici avec le rayon $1$.\footnote{« Sinus rectus » est
le sinus, « sinus versus » le sinus verse $1 - \cos$, « sinus complementi » le
cosinus, « sinus totus » le rayon.} La hauteur d'un point du colure au-dessus
du plan de l'équateur est le sinus de sa déclinaison. Les étapes, avec les
nombres du feuillet :

\begin{enumerate}
  \item[1.] $A$ est à $\alpha A = \varepsilon + \beta = 96^\circ 23'$ de
  l'équateur, donc à la hauteur $AB = \sin 96^\circ 23' = \sin 83^\circ 37'
  = 0{,}99380$.
  \item[2.] $\kappa$ est à $\kappa\beta = \beta - \varepsilon = 49^\circ 20'$ de
  l'équateur, à la hauteur $\kappa\lambda = \sin 49^\circ 20' = 0{,}75851$. La
  corde $A\kappa$ descend donc de $AC = 0{,}99380 - 0{,}75851 = 0{,}23529$, et
  le demi-diamètre du petit cercle correspond à une descente
  $AD = 0{,}11764$.
  \item[3.] L'étoile est à l'arc $37^\circ 34' 20''$ de $A$ sur le petit
  cercle ; sa projection $F$ est à la distance $AF = 1 - \cos 37^\circ 34' 20''
  = 0{,}20741$ de $A$, mesurée en rayons du petit cercle.
  \item[4.] Par les triangles semblables, $AG : AD = AF : AE$ avec $AG = 1$ :
  la descente de $F$ au-dessous de $A$ est $AE = 0{,}11764 \times 0{,}20741
  = 0{,}02440$.
  \item[5.] La hauteur de l'étoile au-dessus de l'équateur est
  $EB = 0{,}99380 - 0{,}02440 = 0{,}96940 = \sin\delta$, d'où
  $\delta = 75^\circ 47' 20''$.
\end{enumerate}

Tous ces nombres sont justes : les sinus aux cinq décimales, les deux
soustractions, la moitié, le produit.\footnote{Les valeurs exactes sont
$\sin 83^\circ 37' = 0{,}993800$, $\sin 49^\circ 20' = 0{,}758514$,
$1 - \cos 37^\circ 34' 20'' = 0{,}207415$, et
$\arcsin 0{,}96940 = 75^\circ 47' 22''$. La vérification est de cette lecture ;
la transcription avait déjà vérifié les soustractions et le produit. Dans
« triangle rectangle $ADG$ », la seconde lettre est repassée et peut être un
$B$.} Et la construction est exactement la formule d'aujourd'hui. En effet
$AB = \sin(\beta + \varepsilon)$, $\kappa\lambda = \sin(\beta - \varepsilon)$,
donc
\[
AD = \tfrac12\bigl(\sin(\beta+\varepsilon) - \sin(\beta-\varepsilon)\bigr)
= \cos\beta \,\sin\varepsilon ,
\]
et $AF = 1 - \cos(\lambda - 90^\circ) = 1 - \sin\lambda$. Il vient
\[
\sin\delta = \sin(\beta+\varepsilon) - (1 - \sin\lambda)\cos\beta\,\sin\varepsilon
= \sin\beta\,\cos\varepsilon + \cos\beta\,\sin\varepsilon\,\sin\lambda ,
\]
qui est la formule de passage des coordonnées écliptiques à la déclinaison.
Calculée directement avec les données du feuillet, elle donne
$\delta = 75^\circ 47' 21''$ : l'écart avec la valeur du feuillet est d'une
seconde.

\subsection*{La hauteur du pôle}

Au passage inférieur, l'étoile est au-dessous du pôle d'un arc égal à sa
distance polaire $90^\circ - \delta$ ; la hauteur du pôle est donc
\[
\varphi = h + (90^\circ - \delta).
\]
Avec $\delta = 75^\circ 47' 20''$, la distance polaire vaut
$90^\circ - 75^\circ 47' 20'' = 14^\circ 12' 40''$, et
\[
\varphi = 34^\circ 35' + 14^\circ 12' 40'' = 48^\circ 47' 40''.
\]
Le feuillet écrit $14^\circ 12' 35''$ pour la distance polaire, et par suite
$48^\circ 47' 35''$ pour la hauteur du pôle, « c'est-à-dire de Paris
(Lutetiæ) ».\footnote{L'erreur est une soustraction : $90^\circ$ moins
$75^\circ 47' 20''$ font $14^\circ 12' 40''$, non $14^\circ 12' 35''$. Elle est
répétée à l'identique dans la solution à la règle et dans la solution par les
nombres, et la somme $48^\circ 47' 35''$ est celle du nombre faux ; la
transcription l'avait relevée. Avec la déclinaison exacte, on trouverait
$48^\circ 47' 39''$. Avant « id est Lutetiæ », un « 48 » est biffé.} La valeur
juste, avec les propres données du feuillet, est $48^\circ 47' 40''$. La corde
$s\rho$ de la figure, menée à $\varphi$ du pôle, représente l'horizon.

\subsection*{Les autres observations de l'hiver 1620--1621}

Au verso, sous le titre « Autres observations de la même année », deux autres
méthodes.

\emph{L'étoile polaire.} Le 2 janvier, vers 5 h 30, sa hauteur minimale
est de $46^\circ 2' 30''$ ; le 4 janvier, sa hauteur maximale au méridien au-dessus
du pôle est de $51^\circ 28' 45''$ ; le 5 janvier, le matin, la hauteur minimale
est notée $46^\circ 2' 40''$, et le soir la hauteur maximale
$51^\circ 28' 45''$.\footnote{Plusieurs mots de ces trois lignes n'ont pas été
lus : celui qui suit « h. » aux deux premières, celui qui suit « 5 h 30 »
à la deuxième, et le dernier mot de la première.} Une étoile circumpolaire
passe au méridien au-dessus du pôle à la hauteur $\varphi + p$ et au-dessous à
la hauteur $\varphi - p$, où $p$ est sa distance polaire ; la demi-somme donne
la hauteur du pôle, la demi-différence la distance polaire :
\[
\varphi = \tfrac12\bigl(51^\circ 28' 45'' + 46^\circ 2' 30''\bigr) = 48^\circ 45' 37\tfrac12'',
\qquad
p = \tfrac12\bigl(51^\circ 28' 45'' - 46^\circ 2' 30''\bigr) = 2^\circ 43' 7\tfrac12''.
\]
Ce sont les nombres du feuillet, et ils sont justes.\footnote{Le feuillet
calcule avec la hauteur minimale du 2 janvier. Avec le $46^\circ 2' 40''$ du 5
janvier, tel qu'il est écrit, on aurait $48^\circ 45' 42\tfrac12''$ et
$2^\circ 43' 2\tfrac12''$. Le $7\tfrac12$ de la distance porte trois traits
d'unité sur le feuillet. La fin de la phrase, après « la même que », n'a pas
été lue, sauf un nom qui paraît être celui de Tycho ; « consequens » est une
lecture douteuse.} La méthode a l'avantage de ne rien demander aux
catalogues : elle ne suppose connue ni la position de l'étoile ni
l'obliquité.

\emph{Le Soleil au solstice d'hiver.} Le 23 décembre 1620, la hauteur méridienne
du Soleil est de $17^\circ 45'$ ; on ajoute la parallaxe, $2' 58''$, ce qui
donne $17^\circ 47' 58''$, puis la déclinaison du Soleil, $23^\circ 30' 22''$
(australe, le Soleil étant à la longitude $272^\circ 13'$, soit $2^\circ 13'$
du Capricorne) : la hauteur de l'équateur est $41^\circ 18' 20''$. Le 25
décembre, la hauteur est de $17^\circ 48'$, soit $17^\circ 50' 58''$ avec la
parallaxe, et avec la déclinaison $23^\circ 27' 24''$ (longitude
$274^\circ 16'$) la hauteur de l'équateur est $41^\circ 18' 22''$, « plus
grande à cause des réfractions dont on n'a pas tenu compte ».\footnote{Les
additions sont justes. Les déclinaisons se vérifient par $\sin\delta =
\sin\varepsilon \sin\lambda$ avec l'obliquité de Tycho : on trouve
$23^\circ 30' 23''$ et $23^\circ 27' 21''$, soit un accord à trois secondes
près. Le mot devant « observatas » est bref et peu sûr ; la phrase peut
vouloir dire que la seconde valeur dépasse la première, ou que toutes deux
dépassent la vraie. Le Soleil est dessiné comme un cercle.} La hauteur du
pôle est le complément de celle de l'équateur,
$\varphi = 90^\circ - 41^\circ 18' 20'' = 48^\circ 41' 40''$ ; le feuillet
s'arrête à la hauteur de l'équateur et ne fait pas cette soustraction.

La réfraction atmosphérique relève les astres ; en la négligeant, on
surestime la hauteur de l'équateur et l'on sous-estime celle du pôle. Le
feuillet a raison sur le sens de l'effet.\footnote{À $17^\circ 45'$ de hauteur,
la réfraction moyenne est d'environ $3'$. La parallaxe de $2' 58''$ ajoutée
par le feuillet est celle de l'astronomie de son temps ; la parallaxe solaire
vraie est d'environ $9''$, et l'obliquité de Tycho dépasse d'environ $2'$ celle
de 1621. Ces trois corrections vont dans le même sens et expliquent l'essentiel
de l'écart entre la valeur solaire et les valeurs stellaires. Les deux
méthodes stellaires, elles, ne corrigent pas la réfraction, qui les abaisserait
d'une minute environ. La latitude de Paris est aujourd'hui voisine de
$48^\circ 50'$ ($48^\circ 50' 11''$ à l'Observatoire) ; le feuillet ne dit pas
où les observations furent faites. Ces ordres de grandeur sont de cette
lecture.} Les trois méthodes donnent ainsi $48^\circ 47' 35''$ (corrigé :
$48^\circ 47' 40''$), $48^\circ 45' 37\tfrac12''$ et $48^\circ 41' 40''$ ; le
feuillet ne les compare pas.

\section*{La lettre de Mersenne au P. Petau (vues 39 et 40)}

\pagerange{39}{40}
Un petit feuillet oblong, monté seul, d'une écriture petite, rapide et très
liée, dont beaucoup de mots ne se lisent qu'au sens et dont plusieurs ne
se lisent pas. Il est signé « Mar. Mers. M. », Marin Mersenne, Minime ; une
main plus tardive a écrit au-dessous « Cette lettre est du P. Mersenne,
minime ».\footnote{Le dos du feuillet (vue 40) ne porte que l'adresse, « Au
R. […] père Petau », un mot bref n'ayant pas été lu. Aucune date, aucun lieu,
aucune marque postale. Le chiffre de la série courante à l'encre, en haut à
droite du feuillet, se lit 25 ou 26 ; le catalogue en ligne de la
correspondance de Mersenne cite cette lettre comme « nouv. acq. fr. 9544,
f. 25r ».} Voici ce qui s'en lit, en français d'aujourd'hui ; les crochets
marquent les mots non lus.

Mon révérend père, me souvenant de ce que vous me demandâtes l'hiver passé
sur la hauteur du pôle, pour la latitude de Paris, il m'est souvenu de
quelques observations […] de janvier 1621, […] lesquelles je vous envoie, si
tant est que vous […] ; si vous ne les avez, je vous pourrais servir assez à
propos. Je n'en ai point retenu de copie : c'est pourquoi je vous prie de me
renvoyer celle-ci quand vous vous en serez servi. Je vous prie aussi de
donner ce petit mot de lettre au révérend père Grandamy […] de ce qui est de
la longitude de Paris, s'il vous vient quelque chose de […]. Cependant je me
recommande à vos saintes prières. Votre très affectionné […], Marin Mersenne,
minime.\footnote{Sont donnés avec doute, dans l'ordre : « Révérend »,
« hiver », « passé », « sur », « Jan. », « suivant », « ou », « pourrais »,
« servir », « retenu », « Grandamy », « vous vient », « cependant »,
« saintes », « Vostre ». Le millésime 1621 est net. Au-delà de « quelque chose
de », la ligne porte un mot de six ou sept lettres puis un nombre de quatre
signes, que la transcription n'a pu fixer. Le cachet de la bibliothèque couvre
la fin de la cinquième ligne.}

La lettre et la feuille se répondent : la lettre envoie des observations de
janvier 1621 sur la hauteur du pôle, la feuille les contient et porte au dos
l'adresse de Petau au Collège de Clermont. La lettre ne se date pas
elle-même : elle est postérieure à l'hiver où Petau posa sa question, et aux
observations qu'elle envoie. La seconde demande, qui touche la longitude de
Paris, ne trouve rien sur la feuille : la hauteur du pôle donne la latitude,
non la longitude, qui demandait des observations simultanées en deux lieux,
d'éclipses par exemple.\footnote{La remarque sur la longitude est de cette
lecture ; le passage de la lettre qui la concerne n'est lu qu'en partie.}

\section*{La description des manuscrits de Roberval par Libri (vues 41 à 55)}

\pagerange{41}{55}
Une fiche moderne (vue 41) donne le titre : « G. Libri. Description des
manuscrits de Roberval conservés à la Bibliothèque nationale ». Suivent onze
feuillets d'un format plus grand que l'album, écrits recto et verso d'une
seule main ronde et régulière, que le rédacteur pagine de 1 à 20 ; la
première page porte « Manuscrit de Roberval. par G. Libri ».\footnote{« par
G. Libri » est d'une plume plus fine, peut-être d'une autre main. Rien ne dit
si la pièce est l'autographe de Libri ou une copie. Au pied de la première
page : « Vente B. Boncompagni, VI. Autogr. (1898), n\textsuperscript{o} 673. »
Sur les versos, la fin des lignes passe sous le montage, au pli ; la plupart
des mots perdus du texte le sont là, et non par défaut de lecture. Les nombres
de la marge renvoient aux pages ou aux feuillets des manuscrits décrits.} Le
texte décrit cinq volumes l'un après l'autre, chacun sous son titre courant.
Ce n'est pas un texte de mathématiques mais un inventaire commenté ; on en
donne ici le contenu, et l'on énonce en termes d'aujourd'hui ce qui peut
l'être.

\subsection*{Le manuscrit français au numéro laissé en blanc (vues 42 à 49)}

Un in-folio de 948 pages non numérotées, dont le numéro est resté en blanc
dans la description. Il contient cinq écrits : des éléments de géométrie
(pages 11 à 681) ; un petit écrit sur la quadrature de la parabole, composé
en mars 1651 (685 à 703) ; des propositions sur la comparaison et la mesure
des surfaces (705 à 747) ; des éléments de musique en latin (751 à 818) ; des
observations sur la composition des mouvements et sur le moyen de trouver les
tangentes des courbes (821 à la fin).\footnote{Le feuillet écrit « de la page
821 à la page 643 », lu nettement : lapsus du rédacteur, la dernière page
devant dépasser 821 ; la description suivante mène l'écrit au moins jusqu'à
la page 945. La page 11 porte, selon Libri, le n\textsuperscript{o} 141, et le
volume les lettres « D. d », la seconde douteuse.}

\emph{Les éléments de géométrie.} Roberval prévoit dix livres pour la
géométrie de la pratique ordinaire — cinq pour la géométrie plane, les
rapports et les proportions, trois pour la géométrie dans l'espace, deux pour
les nombres — et promet dix autres livres de spéculation.\footnote{« dix » est
un mot repassé et empâté ; le total cinq, trois et deux le confirme.} Il donne
d'abord les définitions, axiomes, demandes et énoncés des neuf premiers
livres, et quatre définitions pour un treizième livre qui devait reprendre
les livres XIII et XIV d'Euclide.\footnote{Le second ordinal est caché au pli ;
la lecture « XIV » est celle du sens.} Libri critique plusieurs définitions :
celle de l'unité, « ce par quoi une chose est dite une », qui n'apprend rien
à qui ne le sait déjà ; celle du solide, longue et obscure. Il approuve
l'ordre qui définit ensuite la surface comme l'extrémité du solide, la ligne
comme l'extrémité d'une surface et le point comme l'extrémité d'une ligne ;
la droite est la ligne qui reste immobile quand tout tourne autour de deux de
ses points. Roberval nomme « portion de cercle » le segment, « angle de
contingence » l'angle d'une tangente avec le cercle, et définit la tangente
comme une droite passant par deux points consécutifs de la courbe — définition
que Libri rapproche de la théorie des infiniment petits.\footnote{La définition
est soulignée mot à mot dans le manuscrit. Le rapprochement est celui de
Libri : en termes d'aujourd'hui, la tangente est la limite des sécantes quand
les deux points se rapprochent, ce que « deux points consécutifs » dit sans
la notion de limite.} Les triangles sont dits rectangles, amblygones
(obtusangles) et oxygones (acutangles) ; les solides semblables, ceux qui
sont compris sous un même nombre de plans semblables ; la « portion de
sphère » est un segment sphérique, le « trapèze sphérique » la portion comprise
entre deux plans. Libri juge que les définitions n'ont pas toute la précision
voulue.

Chaque proposition porte en marge la proposition d'Euclide dont elle relève,
ou la mention qu'elle est nouvelle. Libri compte les propositions : 32 au
livre I, dont sept nouvelles ; 32 au livre II ; 32 au livre III ; 35 au
livre IV ; 39 au livre V, dont 25 tirées d'Euclide ; 33 au livre VI, dont
deux seulement tirées d'Euclide ; 38 au livre VII, qui est le livre XII
d'Euclide amplifié et réordonné, avec quatre propositions nouvelles ; rien
pour le livre VIII ; 35 propositions élémentaires sur les nombres au livre
IX, sans les irrationnelles d'Euclide.\footnote{Plusieurs nombres, cachés
au pli, manquent : combien de propositions « nouvelles » au livre II, combien
« sans indication » au livre III, combien « tirées d'Euclide » au livre IV.
Le renvoi marginal du livre VII se lit 557, le dernier chiffre pouvant être un
7 ou un 9.} Il relève quelques énoncés, tous justes : il existe des plans, et
l'on peut leur mener une perpendiculaire, une droite tournant autour d'une
droite qui lui est perpendiculaire engendrant un plan ; trois points non
alignés déterminent un plan ; un cercle a un seul centre ; par trois points
non alignés passe un seul cercle ; deux cercles concentriques ne peuvent se
couper sans se confondre ; la position relative de deux cercles se lit sur
la distance $d$ des centres comparée aux rayons $r$ et $r'$ — ils se coupent
si et seulement si $|r - r'| < d < r + r'$ — et réciproquement.\footnote{Pour
les trois points et le cercle, le texte dit « on peut faire passer qu'une
circonférence » : un « ne » a pu disparaître au pli, et l'énoncé correct est
celui de l'unicité. La condition sur $d$ est l'énoncé moderne, que Libri ne
donne pas ; il dit seulement que la proposition établit « les rapports » de la
ligne des centres et des rayons selon la position des cercles.}

Parmi les nouveautés du livre II, Libri cite la septième : des droites
parallèles ont des perpendiculaires communes, démontré à partir d'une
définition des parallèles comme droites équidistantes. Il faut noter ce que
cette définition contient : poser qu'il existe des droites équidistantes
d'une droite donnée équivaut au cinquième postulat d'Euclide, de sorte que la
proposition, juste, repose sur un postulat des parallèles caché dans la
définition.\footnote{Cette remarque est de cette lecture ; Libri n'en dit
rien. Elle est au cœur de la longue histoire des tentatives de démonstration
du postulat des parallèles, et elle n'enlève rien à la validité de la
proposition sous la définition adoptée.}

\emph{La quadrature de la parabole (1651).} Par la méthode des indivisibles,
Roberval donne l'aire du segment parabolique, d'abord du côté convexe, puis
du côté concave ; puis il montre que le conoïde parabolique (le paraboloïde de
révolution) est la moitié du cylindre de même base et de même hauteur, et que
le fuseau parabolique — le solide engendré par un segment de parabole
tournant autour de sa base — est les $8/15$ du cylindre qui l'enveloppe. Ces
deux rapports sont justes.\footnote{En termes d'aujourd'hui : en faisant
tourner autour de l'axe $Ox$ la région $0 \le y \le \sqrt{x}$, $0 \le x \le
a$, on obtient $\pi \int_0^a x\,dx = \tfrac12 \pi a^2$, moitié du cylindre
$\pi a \cdot a$. Pour le fuseau, avec $y = h(1 - x^2/b^2)$ tournant autour de
$Ox$ entre $-b$ et $b$, $\pi \int_{-b}^{b} h^2 (1 - x^2/b^2)^2\,dx =
\tfrac{8}{15}\,\pi h^2 \cdot 2b$. Les calculs sont de cette lecture.} Il
donne ensuite la quadrature de « toutes les paraboles à l'infini », le
triangle étant la parabole du premier degré : en termes d'aujourd'hui,
\[
\int_0^a x^n\,dx = \frac{a^{n+1}}{n+1},
\]
l'aire sous la parabole de degré $n$ est la fraction $1/(n+1)$ du rectangle
circonscrit. Il en tire les rapports des solides engendrés par ces paraboles
tournant autour de leur ordonnée commune, puis autour de leur axe — pour la
courbe $x = y^n$ tournant autour de son axe, le conoïde est au cylindre comme
$n$ est à $n + 2$, ce qui redonne la moitié pour $n = 2$ — puis autour d'un
diamètre quelconque.\footnote{Libri dit que Roberval donne ces rapports sans
les écrire. Le rapport $n/(n+2)$ est calculé par cette lecture comme exemple
de ce que ces propositions établissent ; il n'est pas sur le feuillet. Le
cachet de la bibliothèque couvre une partie de « cylindre » et un mot après
« le second ».}

\emph{La comparaison des surfaces.} L'écrit porte en marge « Archimède, livre
I de la sphère et du cylindre », et en reprend les grands énoncés, tous justes :
la surface de la sphère est égale à la surface latérale du cylindre circonscrit
($4\pi r^2 = 2\pi r \cdot 2r$) ; la surface latérale d'un cylindre droit est
au cercle de base comme le côté est à la moitié du rayon ($2\pi r h : \pi r^2
= h : \tfrac{r}{2}$) ; celle d'un cône droit est à sa base comme le côté est
au rayon ($\pi r l : \pi r^2 = l : r$) ; la sphère est au cylindre circonscrit
comme 2 est à 3, et de même, dit Libri, pour le secteur.\footnote{Plusieurs
fins de ligne sont cachées au pli ; « pour le secteur » est lu par le sens, le
l de « le » étant tracé comme un f.} Vient ensuite un lemme qui ramène les
démonstrations par les indivisibles à la méthode d'exhaustion des Anciens :
si deux quantités sont telles qu'en augmentant ou en diminuant très peu l'une,
leur rapport devienne dans un cas plus grand, dans l'autre plus petit qu'un
rapport donné, elles sont entre elles dans ce rapport. Tel quel, l'énoncé
n'est vrai que si « très peu » s'entend « aussi peu qu'on veut » : si, pour
tout $\epsilon > 0$, on a $(A - \epsilon)/B < r < (A + \epsilon)/B$, alors
$A/B = r$.\footnote{La précision est de cette lecture ; Libri rapporte le
lemme en ces termes et Roberval l'entendait sans doute ainsi. Libri ajoute que
Roberval appuie la proposition de quelques exemples.}

Suivent : l'ellipse est au cercle de même diamètre comme la perpendiculaire
abaissée de l'extrémité du diamètre conjugué sur le premier diamètre est au
rayon du cercle, ce qui est juste, l'aire de l'ellipse de demi-diamètres
conjugués $a'$, $b'$ faisant l'angle $\theta$ étant $\pi a' b' \sin\theta$ ; la
comparaison des conoïdes hyperboliques et du sphéroïde avec leurs cylindres et
leurs cônes ; un article sur les spirales, comparées à leur cercle pour en
conclure l'aire ; deux autres démonstrations du rapport du conoïde hyperbolique
au cylindre ; enfin les conchoïdes, celle de Nicomède d'abord, avec des
comparaisons d'aires, puis des conchoïdes en nombre infini.

\emph{Les éléments de musique.} Un traité latin, \emph{Elementa Musicae}, en six
chapitres : les rapports et proportions de l'harmonie ; le son et ses
accidents, et les intervalles ; les trois genres ; l'espèce, le système et le
mode ; un chapitre très court sur le chant, qui n'examine aucun instrument
malgré son titre ; la durée, la mesure et la valeur des notes, et les
instruments.\footnote{Le titre latin du cinquième chapitre commence au pli,
où un mot est caché. Les trois genres sont, dans la tradition grecque que le
traité suit, le diatonique, le chromatique et l'enharmonique ; l'explication
est de cette lecture.}

\emph{La composition des mouvements et les tangentes.} C'est la méthode des
tangentes de Roberval : une courbe étant décrite par un point animé de deux
mouvements connus, la tangente en un point est la direction de la vitesse
résultante, obtenue en composant les deux vitesses par le parallélogramme.
Libri note que l'écrit est imprimé dans le tome VI des anciens Mémoires de
l'Académie des sciences, mais que l'imprimé s'arrête au treizième exemple, la
tangente à la parabole de Descartes, alors que le manuscrit en contient un
quatorzième, sur une courbe qu'il appelle « Téroïde ou Aile de
Rob[erval] ».\footnote{Les deux mots sont lus lettre à lettre et restent
douteux. Si le premier est « Ptéroïde », du grec \emph{pteron}, aile, les deux
noms disent la même chose ; cette lecture ne l'affirme pas. La « parabole de
Descartes » est le nom de la cubique de la \emph{Géométrie}, aussi appelée
trident.} Suivent cinq planches de vingt-six figures ; une série de
propositions attribuée à « M. de Rob », où l'auteur donne une parabole égale à
une hélice ; une septième proposition qui applique la méthode des mouvements
composés à plusieurs problèmes et retrouve les conséquences que Fermat avait
tirées de sa méthode \emph{de maximis et minimis} ; enfin une copie de lettre,
sans signature ni adresse, sur les oscillations de corps de figures
diverses.\footnote{« Hélice » désigne au XVII\textsuperscript{e}
siècle aussi la spirale d'Archimède ; si c'est le cas ici, la proposition est
l'égalité, bien connue alors, d'un arc de spirale et d'un arc de parabole.
Libri ne précise pas. Le dernier chiffre du renvoi « 926 » est repassé.}

\subsection*{Le manuscrit latin 7226 (vues 51 et 52)}

Un in-folio relié de 532 pages non numérotées, qui contient : un
\emph{Tractatus mechanicus} de Roberval daté de 1648 (pages 7 à 70) ; une
lettre de Roberval à Fermat, « conseiller de Toulouse » (71 à 111) ; une
proposition pour trouver les centres de gravité (112 à 116), adressée à Fermat
le 1\textsuperscript{er} avril 1645 et réduite à des définitions et des
énoncés ; un \emph{Theorema lemmaticum} sur les centres de gravité, daté de
1645 (121 à 167) ; un traité de mécanique « et spécialement de la conduite et
élévation des eaux » (173 à 421).\footnote{Le nom de Fermat est écrit
« Fermates » dans la description, finale nettement formée. Le mot latin lu
« nisu » dans le titre du \emph{Theorema lemmaticum} n'est pas sûr. La page 7
porte, selon Libri, le n\textsuperscript{o} 289 et la lettre A.}

Le \emph{Tractatus mechanicus} pose des postulats et des définitions, traite
de l'équilibre de forces situées dans un même plan, définit un « centre de la
vertu de la puissance », et se termine par la démonstration mécanique de la
loi du levier : deux poids en raison inverse des bras se font équilibre. La
lettre à Fermat démontre de même l'équilibre du levier et envoie à Fermat « un
principe de géostatique » qu'il avait demandé. Le \emph{Theorema lemmaticum}
distingue trois cas selon la position de points donnés par rapport à un plan
— tous du même côté ; certains sur le plan, les autres d'un même côté ; des
deux côtés, avec parfois des points sur le plan — et l'applique au
demi-cercle, à la demi-circonférence, à la trochoïde, à sa « compagne » et au
triangle.\footnote{La trochoïde est la cycloïde ; sa compagne, la courbe que
Roberval associait à la cycloïde pour en quarrer l'aire, est une sinusoïde.
Si le lemme porte, comme les trois cas le suggèrent, sur la distance du centre
de gravité d'un système de points à un plan, sa forme moderne ne fait qu'un
cas : la distance algébrique du centre de gravité au plan est la moyenne
pondérée des distances algébriques des points. C'est une interprétation de
cette lecture ; Libri ne donne pas l'énoncé. « à la demi-circonférence » lit
« à la d[…] circonférence », le mot étant caché au pli.}

Le traité des eaux explique d'abord les machines simples — balance, levier,
roue et essieu, poulies et moufles, plan incliné, coin, vis, marteau — puis
la pesanteur des corps plongés dans les liquides, les formes propres à nager,
la conduite naturelle des eaux (siphon, écluses), et la conduite forcée, dont
Roberval dit qu'elle se fait de deux manières, par une pression plus forte
que la pesanteur et par « la crainte du vide » : contrepoids, moulins,
chapelets, vis d'Archimède, « escalier », pompes aspirantes et
foulantes.\footnote{La « crainte du vide » est l'\emph{horror vacui} de la
physique ancienne, que la pression atmosphérique a remplacée comme
explication de l'aspiration ; le mot est celui du manuscrit, rapporté par
Libri.} Le volume se termine par le principe d'Archimède, que Roberval fonde
sur l'action de l'eau qui soulève le corps par-dessous ; d'où il conclut qu'un
corps même plus léger que l'eau, appliqué contre le fond sans liquide
au-dessous de lui, ne remonterait pas. La conclusion est juste : la poussée
est la résultante des pressions sur la surface mouillée, et si la face
inférieure ne l'est pas, la résultante des pressions est dirigée vers le
bas.\footnote{L'explication par la résultante des pressions est de cette
lecture. Le mot « il », avant « ne remonterait », est surchargé.}

\subsection*{Le supplément français 597 (vue 53)}

Un in-folio relié de 430 pages, dont la page 9 porte le
n\textsuperscript{o} 290.\footnote{À la suite du titre courant, d'une autre
encre et semble-t-il d'une autre main, « (Fr. 12279) », les deux derniers
chiffres surchargés : une concordance ajoutée plus tard, que cette lecture ne
vérifie pas.} Il contient un « livre troisième de dioptrique, ou des
lunettes » (pages 7 à 197) ; un traité de Roberval sur la résolution
géométrique des équations planes et cubiques (247 à 309) ; un traité de la
division des lieux en divers degrés (310 à 329) ; un traité pour ramener un
lieu géométrique donné à une équation analytique (330 à 354).

Le premier ouvrage paraît à Libri être du P. Mersenne, dont les deux premiers
livres auraient été publiés à la suite du traité de perspective du P. Niceron,
et ce troisième livre, incomplet, être l'analyse d'un traité d'optique de
Roberval.\footnote{Le nom est coupé au pli après « Mer ». L'attribution est
celle de Libri.} Il traite de la réfraction et de la densité des corps, du
rapport de réfraction et de sa détermination — l'indice de réfraction —, des
termes de la dioptrique, de la différence entre réfraction et réflexion, des
lunettes simples, dont il distingue dix-huit cas, des lunettes composées d'un
ou de plusieurs verres ; une dix-septième proposition, seulement énoncée,
devait traiter de l'œil, de ses défauts et des lunettes qui les corrigent.
Sept planches, huit figures. Le traité des équations n'est, dit Libri, que le
commencement d'un travail plus étendu imprimé dans le tome VI des anciens
Mémoires.

Le traité des lieux reprend la classification des Anciens : il y a une
infinité de degrés de lieux ; au premier genre, les lieux qui sont des lignes,
au second ceux qui sont des surfaces. Les lignes se répartissent en lieux
plans (la droite et le cercle), lieux solides (les trois coniques) et lieux
linéaires (toutes les autres courbes), et l'auteur ne retient que ceux qui
s'enferment dans une équation. En termes d'aujourd'hui, les lieux plans sont
les courbes de degré 1 et le cercle, les lieux solides les courbes de degré 2,
c'est-à-dire les coniques ; la classification est celle de Pappus, et la
restriction aux courbes qui ont une équation est le programme de la géométrie
analytique. Le dernier écrit l'applique au cercle, à la parabole, à
l'hyperbole, à l'ellipse et à la conchoïde.\footnote{Le renvoi « 332 » a un
dernier chiffre incertain, 2 ou 7. Libri dit ces deux derniers ouvrages
« imprimés avec le précédent ».}

\subsection*{Le supplément latin 218, et ce qu'il est devenu (vues 54 et 55)}

Le supplément latin 218, dit Libri, se compose de trois volumes in-quarto
reliés.

\emph{Le tome I} porte au recto de son deuxième feuillet le
n\textsuperscript{o} 171 et le titre \emph{De recognitione aequationum}, au
verso le n\textsuperscript{o} 26 ; l'ouvrage commence au feuillet suivant et
compte 37 pages.\footnote{Le premier mot du titre latin est caché au pli ;
la lecture « De » est celle du sens.} Libri en donne le plan : les équations
du second degré en quatre chapitres (deux racines positives ; deux négatives ;
une de chaque signe ; l'équation qui « subsiste par elle-même ») ; les
équations cubiques en six chapitres, suivis de leur « détermination » ; les
équations du quatrième degré en dix chapitres, jusqu'aux racines
« fictives ». Il ajoute qu'un ouvrage du même titre est imprimé dans le tome
VI des anciens Mémoires de l'Académie des sciences, mais que la rédaction et
la disposition des matières ne ressemblent pas à celles du manuscrit. Le
volume se termine, dit-il, par une petite dissertation latine sur l'art de
voler, et par un article italien avec figures, où l'auteur soutient qu'il
n'est pas impossible à l'homme de voler et fait intervenir les noms de Galilée
et d'Archimède.\footnote{En termes d'aujourd'hui, les racines « négatives »
sont les racines réelles négatives, les « fictives » les racines complexes non
réelles ; la « détermination » est l'étude des racines multiples.}

Ce tome I est l'actuel Latin 11195. La transcription de ce volume, faite dans
ce projet, relève sur les feuillets exactement ce que Libri décrit : l'ancienne
cote « Suppl. lat. 218 » sur la garde et sur la première page du traité ; au
recto du deuxième feuillet (vue 2), le titre \emph{De Recognitione Æquationum}
sous un « 171 » à l'encre ; au verso (vue 3, page de gauche), un « 26 » ; le
traité commençant à la page suivante, paginé par son scripteur de 1 à 37 et
s'achevant à la vue 23 ; quatre chapitres pour les équations du second degré,
six pour les cubiques, dix pour les quarto-quadratiques ; puis le dossier du
vol, avec un feuillet de titre « Ars Volandi », des notes latines, et un traité
italien illustré de deux dessins, qui nomme Galilée, Sagredo et
Archimède.\footnote{Voir \emph{transcripts/latin-11195/batch-01.fr.tex}, en-tête
et vues 1 à 3, et \emph{batch-02.fr.tex}, en-tête. La garde de la vue 1 porte
« Suppl.\textsuperscript{t} lat. 218. A. », la page de la vue 3
« Suppl.\textsuperscript{t} Lat. n\textsuperscript{o} 218. » et
« Suppl.\textsuperscript{t} l. 218. + I. ». La « petite dissertation latine » de Libri correspond aux notes latines des vues 26 à 29, qui
comprennent la page de titre recopiée de l'\emph{Ars volandi} de Flayder.}
Une différence doit être dite. Libri range le traité sur les équations parmi
les manuscrits de Roberval ; les feuillets du Latin 11195 ne nomment Roberval
que comme l'auteur, « D. Rob. », de lemmes que le traité emprunte, et la
lecture modernisée de ce volume n'attribue le traité à personne. Le classement
de Libri est un fait sur Libri ; il n'est pas une preuve.

\emph{Le tome II}, de 41 feuillets, porte le n\textsuperscript{o} 27 au verso
du feuillet 2 et le n\textsuperscript{o} 119 en haut du feuillet 3. Il
contient un petit écrit d'optique, \emph{Appendix ad tractatum de legitimo
Dioptrarum usu}, de cinq propositions (deux théorèmes, trois problèmes) avec
huit planches et vingt-quatre figures, qui renvoie souvent au traité des
coniques de son auteur ; un problème adressé au P. Mersenne le 10 novembre
1642, « trouver le cylindre de plus grand \emph{ambitus} inscrit dans une
sphère donnée » ; l'extrait d'une lettre au P. Mersenne du 3 juin 1639 qui
annonce de nouvelles hélices avec leur démonstration ; et le \emph{De Vacuo}
de Roberval, adressé à Des Noyers aux ides de mai 1648, qui rapporte ses
propres expériences sur le vide, celle de Pascal sur la cause de l'ascension
du mercure dans le baromètre, et des expériences semblables faites sur les
petites montagnes des environs de Paris.\footnote{La parenthèse qui suit le
renvoi aux coniques est lue lettre à lettre (« mo. x. lib. nothorum
conicorum ») et reste douteuse. « Invenire » est écrit avec une initiale en
forme de S. Les ides de mai sont le 15 mai.}

Le problème de 1642 se résout aisément aujourd'hui, et sa réponse dépend du
sens d'\emph{ambitus}. Soit $R$ le rayon de la sphère, $r = R\cos t$ le rayon
du cylindre inscrit et $2R\sin t$ sa hauteur. Si l'\emph{ambitus} est la
surface latérale, $4\pi R^2 \sin t \cos t = 2\pi R^2 \sin 2t$ est maximale pour
$t = 45^\circ$, et vaut $2\pi R^2$. Si c'est la surface totale,
$2\pi R^2(\cos^2 t + \sin 2t)$ est maximale pour $\tan 2t = 2$, et vaut
$\pi R^2 (1 + \sqrt5)$.\footnote{Ces deux solutions sont de cette lecture. Libri
ne donne que l'énoncé, et le feuillet ne dit pas quelle surface Roberval
entendait.}

\emph{Le tome III}, de 68 feuillets, porte le n\textsuperscript{o} 28 au verso
du feuillet 1 et le n\textsuperscript{o} 172 au recto du feuillet 2. Libri y
compte dix lettres, dont certaines sont des écrits : une lettre de Roberval à
Torricelli datée de Paris aux calendes de janvier 1646 (le 1\textsuperscript{er}
janvier), où il revendique la priorité des découvertes sur la cycloïde et
attribue à la France presque toutes les inventions faites sur cette courbe ;
une lettre de Huygens, datée de Leyde en octobre 1646, sur le mouvement des
corps, où l'auteur soutiendrait que la parabole décrite par les projectiles
n'est pas un effet de l'attraction terrestre ; une lettre de Torricelli à
Roberval, de Florence, 7 juillet 1646, avec plusieurs problèmes sur la
cycloïde ; une lettre à Roberval contenant un théorème sur une sorte de
spirale ; une lettre de Torricelli au P. Mersenne, de la même date ; un
théorème sur l'hyperbole ; une lettre de Torricelli à Carcavi, 8 juillet 1646,
sur un problème de Fermat sur les nombres et sur la priorité de la cycloïde ;
une lettre de Roberval à Torricelli, imprimée dans le tome VI des anciens
Mémoires ; un « nouvel usage en analyse des radicaux du second degré et des
degrés supérieurs » par Fermat, avec un appendice ; enfin la copie d'une
lettre de Fermat à Carcavi du 20 août 1650 sur le « débrouillement des
asymétries ».\footnote{Le nombre de feuillets, « 68 », a son second chiffre
empâté. Le quantième de la lettre de Huygens est caché au pli : seul se voit
un premier chiffre 2. « Carcavi » se lit « Careavi » à la lettre. Libri donne
la même date, 7 juillet 1646, aux deux lettres de Torricelli. Les
« asymétries » sont les quantités irrationnelles, et leur « débrouillement »
l'élimination des radicaux d'une équation ; la glose est de cette lecture.
Libri résume la lettre de Huygens en une phrase, et cette lecture ne la
corrige pas faute d'avoir le texte.}

Ces deux derniers volumes sont, sous leurs cotes d'aujourd'hui, les Latin
11196 et 11197, que la notice de la bibliothèque groupe avec le Latin 11195
sans en décrire le contenu, et qui ne sont pas numérisés : la description de Libri est ici le seul inventaire de
leur contenu que le projet ait sous les yeux. Laquelle des deux cotes répond
à quel tome, les feuillets de cet album ne le disent pas.\footnote{Le texte
s'arrête sur un trait horizontal, à la vue 55, qui clôt la description du
supplément latin 218.}

\end{document}
